% AI prompts of the project (course requirement), built into AI_PROMPTS.pdf:
%   pdflatex AI_PROMPTS.tex   (run twice)
\documentclass[11pt,a4paper]{article}
\usepackage[T1]{fontenc}
\usepackage{newtxtext}
\usepackage[a4paper,margin=2.1cm]{geometry}
\usepackage{xcolor}
\usepackage{tikz}
\usepackage[most]{tcolorbox}
\usepackage{booktabs}
\usepackage{tabularx}
\usepackage{titlesec}
\usepackage[hidelinks]{hyperref}
\usepackage{url}

\setlength{\parindent}{0pt}
\setlength{\parskip}{4pt}
\urlstyle{tt}

\definecolor{cnaman}{HTML}{1F5FAD}
\definecolor{cvivek}{HTML}{B35C00}
\definecolor{canshu}{HTML}{1E7B4F}
\definecolor{castitva}{HTML}{7A3EA8}

% name tag of the team member whose part of the work the prompt belongs to
\newcommand{\tagbox}[2]{\tikz[baseline=(n.base)]\node[draw=#1,fill=#1!10,text=#1!80!black,rounded corners=2pt,
  inner xsep=4pt,inner ysep=1.6pt,font=\sffamily\bfseries\small](n){#2};}
\newcommand{\Naman}{\tagbox{cnaman}{Naman}}
\newcommand{\Vivek}{\tagbox{cvivek}{Vivek}}
\newcommand{\Anshu}{\tagbox{canshu}{Anshu}}
\newcommand{\Astitva}{\tagbox{castitva}{Astitva}}

% one prompt: number, date, member tag(s)
\newtcolorbox{prompt}[3]{breakable,enhanced,colback=white,colframe=black!25,boxrule=0.5pt,arc=1.5pt,
  colbacktitle=black!5,coltitle=black,toptitle=2pt,bottomtitle=2pt,left=7pt,right=7pt,top=5pt,bottom=5pt,
  before skip=4pt,after skip=9pt,
  title={\textbf{#1}\hspace{0.7em}#3\hfill{\small\color{black!65}#2}}}

\titleformat{\section}{\large\bfseries}{}{0pt}{}
\titleformat{\subsection}{\normalsize\bfseries\color{black!80}}{}{0pt}{}
\titlespacing*{\section}{0pt}{14pt}{6pt}
\titlespacing*{\subsection}{0pt}{10pt}{4pt}

\begin{document}

\begin{center}
{\LARGE\bfseries AI Prompts}\\[5pt]
{\large EE2000 Final Project Report (FPR-8), Team 8}\\[3pt]
Fetal ECG Extraction from Maternal Abdominal ECG and Fetal QRS Detection on FPGA
\end{center}

\vspace{4pt}
This file lists the prompts we gave to AI tools during the project, in time order and grouped by phase.
Part~A covers the Antigravity agent tool (Google), used from the mid-term phase up to the first FPGA implementation
(23 October to 6 November 2026). Part~B covers Claude Code (Anthropic), used in the final phase (7 to 14 November 2026).
The tag next to each prompt names the team member whose part of the work the prompt belongs to, following the division
of work in Sec.~27 of the final report. Dates are in IST.

\begin{center}
\begin{tabularx}{\textwidth}{@{}lX@{}}
\toprule
Member & Part of the work \\
\midrule
\Naman & Data, test set-up and integration: data set preparation, test cases and test suite, web demo, report assembly \\
\Vivek & Fixed-point design and first full FPGA version: fixed-point model, UART and PC--FPGA protocol, v5 on MicroBlaze, board evaluation and board runs \\
\Anshu & Real-time algorithm and accuracy: multichannel study, U-Net baseline, causal algorithm v2--v5, CinC protocol and scoring, v6.1 ideas, look-ahead algorithms v7-LA and v7-DL \\
\Astitva & Final hardware design: hardware/PC split, bit-accurate C model, Verilog RTL, testbenches, board verification, block RAM optimisation, v6.1 RTL and look-ahead unit \\
\bottomrule
\end{tabularx}
\end{center}

% ---------------------------------------------------------------------------------------------
\section{Part A: Antigravity (Google), 23 October to 6 November 2026}
Model: Gemini~3 (Claude Opus~4 for A28).

\subsection{Mid-term phase: repository audit, deep-learning baseline and multichannel study}

\begin{prompt}{A1}{23 Oct 2026}{\Naman}
Conduct a systematic audit of the repository to establish the current project status: completed milestones, pending
deliverables and the immediate next steps in the pipeline. Then execute every validation, simulation and benchmark
script that must pass before the mid-term report deliverables are prepared.
\end{prompt}

\begin{prompt}{A2}{23--24 Oct 2026}{\Anshu}
Locate the training scripts and configuration of the deep-learning model and inspect its architecture and data loaders
without starting a training run. Provide step-by-step instructions to train it on our Windows laptop (ASUS TUF) with full
CUDA acceleration on the NVIDIA GeForce RTX~3050 Mobile GPU: confirm that the virtual environment is active and that
PyTorch detects the GPU, and state whether the data loader traverses the data set folders recursively or needs explicit
paths. During training the composite loss diverged to NaN from epoch 14 (learning rate 0.000206): identify the source of
the numerical instability (exploding or vanishing gradients, division by zero in the loss terms) and add numerical
safeguards. Then analyse and fix the errors raised when retraining with \texttt{-{}-lr 1e-4} and \texttt{-{}-lr 5e-5}.
\end{prompt}

\begin{prompt}{A3}{23--24 Oct 2026}{\Anshu}
Check whether the codebase already contains multi-channel processing scripts or earlier exploratory modules written by
team members, and produce a table comparing multi-channel and single-channel fetal ECG extraction on the standard
metrics. Explain why our reporting has so far emphasised single-channel results, and add an explicit benchmark of the
4-channel configuration used in the hardware and software pipelines: extend the extraction code to full multi-channel
processing, give the command lines that reproduce the numbers, and recompute the combined comparison table. Investigate
the performance drop on Case~5 (physiological noise or electrode artefacts). Finally, compare 4 electrodes with 8, 10 or
16 in terms of FPGA resources, routing complexity and physiological information gained.
\end{prompt}

\begin{prompt}{A4}{24 Oct 2026}{\Anshu}
Training on all target subsets has finished. Describe the verification protocol and the evaluation scripts for
benchmarking the resulting checkpoint. State whether the network was trained on single-channel or 4-channel recordings,
how the training and generalisation-test recordings were chosen, and why the 4-channel deep model is presented as a
software comparison rather than the FPGA datapath despite its higher accuracy. Give a physical and physiological
explanation of the large signal differences between the abdominal electrode channels.
\end{prompt}

\begin{prompt}{A5}{24 Oct 2026}{\Naman}
Review the mid-term report (MTR-8) guidelines together with the LaTeX source of the preliminary report (PPR-8), then
write the complete MTR-8 in LaTeX with all benchmark results, hardware resource estimates and pipeline diagrams. Deepen
it: expand the mathematical formulation, the data tables and the architecture discussion, and cover all 10 simulations at
all 5 SNR levels of the benchmark database rather than only the 5 baseline cases. Fix the tables that overflow the
margins, make sure the compiled PDF is in the distribution package, and explain the Overleaf settings (compiler, main
file) needed to build it. Reorganise the repository into a modular folder structure, with the reports sorted into
\path{preliminary/} and \path{mid_term/}. Finally, write a complete README (repository layout, algorithms, hardware,
reproduction steps), commit it and push.
\end{prompt}

\subsection{Mid-term phase: report review and corrections}

\begin{prompt}{A6}{25--26 Oct 2026}{\Naman}
Review the Team~22 reference report and our improvement-audit document in the report folder, and identify the
architectural, numerical and structural improvements our mid-term report needs. The report currently fails to compile
(log analysis pasted: unescaped underscores inside \texttt{\textbackslash texttt}, a \texttt{verbatim} block inside the
\texttt{\textbackslash newcontent} argument, undefined labels and citation keys, duplicate source files): fix the fatal
errors first, then every reference, citation and table-overrun warning. A local TeX Live distribution is now installed:
compile and check the PDF locally with pdflatex/latexmk, and remove obsolete build files and duplicate report folders.
\end{prompt}

\begin{prompt}{A7}{25 Oct 2026}{\Naman}
Two independent audits of the revised PDF are pasted. Their verdict: the report is much more complete, but not all
issues are resolved, and some sections present planned work or software-model results as stronger evidence than the
validation supports. Address their priorities. Make the evidence defensible: document the exact U-Net split, take every
reported mean from one source of results, and remove unsupported claims of synthesis, RTL verification or board
execution. Repair the experimental comparison: a properly configured reference-channel ANC baseline, stronger
multichannel baselines, and a replication scorecard traceable to published results. Complete the evaluation: confidence
intervals, FHR error analysis, worst cases, causal/adaptive PCA tests and real-data validation. Then use
\path{MTR8_audit_and_fix_plan.md} as the specification and fix every mathematical, numerical and typographical item in
it, in order.
\end{prompt}

\begin{prompt}{A8}{26 Oct 2026}{\Naman}
Apply all corrections in \path{MTR8_v2_audit_python_theory.md}, so that the Python results and the theoretical claims
agree. Move the whole report to the IEEE/CVPR two-column template required by \path{project.pdf}, assemble the mid-term
submission archive as \path{project.pdf} specifies, with a compliance matrix of any missing deliverables, update the
README with the latest results and repository layout, and commit and push everything.
\end{prompt}

\begin{prompt}{A9}{27 Oct 2026}{\Naman}
Use \path{antigravity_prompt_mtr8_corrections.md} (our correction specification, in \path{report/mid_term/}) as the main
directive: audit MTR-8 and apply every correction it lists. Make the report anonymous as the evaluation rules require:
put EE2000 in the course header and remove all student and personal names. Check the DOIs and URLs of references [4],
[5], [6] and [23] and replace broken links with working authoritative sources. Update Section~16 (AI Usage Disclosure) to
describe the debugging done with AI, including the data set normalisation bug and the LaTeX build fixes. Finally, clean
the build files, check that sources, figures and zip agree, build the final submission archive and push.
\end{prompt}

\subsection{Validation campaign and pre-FPGA specification}

\begin{prompt}{A10}{31 Oct 2026}{\Astitva}
Formulate a roadmap for moving the causal fetal ECG pipeline from the Python numerical models to synthesisable Verilog
RTL for the Spartan-7 XC7S50. Diff the active feature branch against main to identify the recent RTL and testbench
changes, run the complete simulation and regression suite to check whether they affect accuracy or timing, and write a
technical reference that summarises the Spartan-7 and Boolean Board datasheets (specifications, pinout, timing
constraints).
\end{prompt}

\begin{prompt}{A11}{31 Oct 2026}{\Naman}
Add a data set characterisation section covering FECGSYNDB, CinC 2013 and NIFEA as stored in \path{data/}. Then run a
complete validation campaign of the real-time multi-channel fetal ECG extraction and QRS/FHR system, acting as
signal-processing, FPGA-verification and test engineer: audit the repository and reconstruct the pipeline that actually
runs (PCA/SVD versus spatial beamforming) while preserving the original; reproduce the existing results; validate on the
full synthetic data set by SNR, physiological condition and subject; test the four-channel configuration (channel
ablation, permutation, channel failure, weight robustness); measure extraction quality and maternal residue; validate
QRS detection and FHR (parameter sensitivity, heart-rate range); run an ablation study; validate on the NIFEA clinical
recordings; check fixed-point overflow and quantisation, real-time streaming, the UART protocol and FPGA resource
readiness; and add automated failure analysis and improvement experiments. All results must be reproducible and end in
one master report, from executive summary to improvement roadmap.
\end{prompt}

\begin{prompt}{A12}{31 Oct 2026}{\Anshu}
Based on the validation results, improve the existing algorithm without rewriting the project or starting RTL. Extend
the 300-case benchmark, which covers only layout l1, to layouts l2--l5 with the same algorithm, beat-matching tolerance,
SNR levels, conditions and subjects, and report the results per layout, subject, SNR and condition. Then investigate
spatial calibration, detector improvements and FHR estimation, look into the NIFEA results and correct the fixed-point
issues, keeping the baseline implementation unchanged.
\end{prompt}

\begin{prompt}{A13}{31 Oct 2026}{\Vivek}
Complete all remaining validation, fixed-point verification and hardware specification work before RTL, for the
RealDigital Boolean Board (Spartan-7 XC7S50-1CSGA324C). Rules: inspect before changing, preserve the baseline, no
fabricated results, no data leakage, everything reproducible. Tasks: reproduce the full-layout benchmark; validate the
autocorrelation-based spatial calibration (selection methods, maternal periodicity failures, heart-rate search range,
confidence and fallback, stability); separate raw QRS performance from FHR post-processing; validate the detector
configuration; verify the fixed-point implementation (beat matching, exact number formats, bit-accurate reference
model); finalise the FPGA specification (module architecture, calibration responsibility, safe weight updates, UART data
protocol, streaming timing, FHR output semantics); estimate resources; document the NIFEA limitations; freeze the
algorithm; and add automated regression tests. Do not start full RTL; deliver one complete ZIP.
\end{prompt}

\begin{prompt}{A14}{31 Oct 2026}{\Vivek}
An independent review of the pre-FPGA package found discrepancies between the reports, CSV files, Python code,
regression tests and hardware specifications. Correct them (a correction phase, not new tuning): reconcile the QRS
metrics; investigate the calibration confidence and the PC3 fallback; correct the fixed-point reference model (filter
state widths, negative-number rounding, runtime coefficient generation, stage-by-stage comparison, independent arithmetic
tests); correct M4, interval fusion and the FHR implementation; resolve the FHR range inconsistencies; reconcile the
benchmark coverage; freeze the exact digital filter specifications; make the regression suite runnable on its own; and
check the UART and hardware specifications for consistency. Then run the Phase~3 numerical verification over all
benchmark recordings to fix bit-accurate baseline criteria before RTL synthesis.
\end{prompt}

\begin{prompt}{A15}{31 Oct 2026}{\Anshu}
Determine how much delay the complete signal-processing pipeline introduces and what correction the detected QRS
timestamps need (no calibration, FHR, fixed-point or RTL work in this task): reproduce the existing timing results,
correct the timing evaluation method, trace the actual sources of delay, choose the compensation and add regression
tests. Then, with 12 samples as the leading value, find out why agreement within $\pm$2 samples exceeds 90\,\% on some
recordings but is below 10\,\% on others: keep the current implementation, plot the worst recordings, find the source of
the variation, improve the evaluation method and compare compensation strategies before deciding.
\end{prompt}

\subsection{UART bring-up on the Boolean Board}

\begin{prompt}{A16}{3--4 Nov 2026}{\Vivek}
Run the standalone Verilog testbench in Vivado/XSim and diagnose the failing assertion from the simulation log. Then
carry out a pre-integration audit of our standalone UART echo design (115200 baud 8N1, 100\,MHz clock), which already
passes physical tests (single characters, strings, 512-byte transfers, a full 256-byte sweep, 50 stress iterations): you
prepare each test, I run it on the board and return the evidence, and you analyse, fix and retest. Cover the required
verification groups, Vivado source management, preparation for FECG integration, documentation and git checkpoints. Run
the regression suite (FIFO, timing and framing testbenches) and the production timing testbench, trace the failing timing
cases to the baud generator or the sampling point, and re-run with the corrected sampling.
\end{prompt}

\begin{prompt}{A17}{4 Nov 2026}{\Vivek}
All five corner-case tests now pass (framing error and recovery, reset during TX and RX, FIFO pointer wraparound): audit
whether each test really exercises the intended hardware and whether any PASS could be a false positive. Add a FIFO
saturation and overflow test (fill all 16 entries, reject the 17th byte, drain and check the order). Verify the RTL with
the real board parameters (100\,MHz, 115200 baud, \path{CLKS_PER_BIT} = 868) in a dedicated production timing testbench
with data tests and per-bit timing assertions. When it reports bit durations of 869 and 867 cycles and an RX midpoint at
cycle 0, find out whether the RTL or the testbench edge counting is wrong, without changing the RTL or relaxing any
assertion; after a git checkpoint, correct the TX bit-duration asymmetry; then find why the production testbench stalls
in Test Group~1 until the watchdog expires, again without touching the RTL or the watchdog.
\end{prompt}

\begin{prompt}{A18}{4 Nov 2026}{\Vivek}
Use Teamwork mode with sub-agents (RX RTL audit, testbench measurement audit, independent timing reviewer) to resolve the
production timing assertion ``RX 2-stage synchronizer latency mismatch: expected 2, got 1'', and resume the interrupted
investigation without repeating finished work. Report the status of the RTL changes, the latest stress, timing, unit and
loopback regressions, the git state and the remaining verification items.
\end{prompt}

\begin{prompt}{A19}{4 Nov 2026}{\Vivek}
The production timing testbench passes in Vivado 2025.2: proceed to synthesis, implementation and bitstream generation
for the Boolean Board (XC7S50-1CSGA324C, 100\,MHz clock, UART 115200 8N1), after parallel audits of the RTL, the XDC
constraints and the Vivado project. Trace the UART and clock connections in the Boolean V2I1 schematic (page~9 USB-UART,
page~11 FPGA pins: \texttt{clk} F14, \path{uart_rxd} U11, \path{uart_txd} V12, not V15/W15), audit the whole XDC (I/O
standards, duplicate constraints, unconstrained ports, LED pins traced from the schematic) before synthesis, and correct
\texttt{led[7:0]}. Document the LED diagnostic map and a board test procedure for UART packets. UART stopped working after
the pin change: investigate against the schematic, rebuild with the corrected constraints, keep the verified pins and
test on the board. Then organise the UART folder (RTL, testbenches, logs, constraints), commit, check that the branch
differs from main only in UART files, and merge it into the causal FPGA pipeline branch.
\end{prompt}

\subsection{Python-to-Verilog conversion and board validation}

\begin{prompt}{A20}{4 Nov 2026}{\Astitva}
Teamwork mode, master task: convert the non-ML Python signal processing into synthesisable Verilog for the Spartan-7.
First audit and freeze the Python reference in \path{src/fecg/} (baseline removal, filtering, QRS detection, FHR;
machine-learning models excluded) with reproducible stage-by-stage vectors (milestone M0); map each block to hardware and
write the fixed-point specification with bit widths, signedness, rounding and saturation (M1); audit the UART modules,
the packet protocol and the known watchdog bug; and define the single-channel architecture, resource estimates and
roadmap, pausing for review before any RTL is written. After approving the task breakdown, start RTL under strict
conditions: keep the Python code and golden vectors, never change the frozen algorithm or fixed-point specification
silently, test every module against the fixed-point Python model with exact matches, one module at a time, single-channel
first with the beamformer kept separate, and no hardware claims from simulation. Begin with \path{biquad_df2t.v} only.
\end{prompt}

\begin{prompt}{A21}{4 Nov 2026}{\Astitva}
Independently audit \path{biquad_df2t.v} rather than trusting the report of 22,404 bit-exact comparisons: read the RTL,
its testbench, the checker and the Python reference, and derive every product, accumulator and state width yourself. Then
audit the source line by line: operand signedness, expression widths before assignment, sign extension, arithmetic
shifts, saturation constants, and rounding against Python. Then implement \path{hp1_baseline.v} from the exact recurrence
in \path{src/fecg/}, documenting widths, scaling, rounding, reset behaviour and latency before writing the RTL.
\end{prompt}

\begin{prompt}{A22}{4 Nov 2026}{\Astitva}
Implement the rest of the single-channel RTL in one coordinated sprint, with common interfaces and fixed-point contracts
fixed first and the modules developed in parallel: \path{hp1_baseline}, \path{biquad_df2t}, \path{lp40_filter},
\path{bandpass_config_c}, \path{squarer_mwi}, \path{adaptive_detector}, \path{peak_localizer}, \path{interval_fusion} and
\path{fhr_calculator}, integrated into a verified single-channel pipeline. Then validate it with the real toolchain:
compile and simulate every module and the top level in Vivado (no Python emulation in place of simulation), compare every
stage with the Python golden reference and report exact mismatch counts, separate the latency components instead of
quoting one number, run synthesis and implementation (DSP, LUT, FF and BRAM use, worst slack, maximum clock), program the
board and compare its UART outputs with Python. Do not start the multichannel PCA yet.
\end{prompt}

\begin{prompt}{A23}{4 Nov 2026}{\Astitva}
Phases 2 and 3: get the single-channel design running on the physical Spartan-7 board, implement the four-channel PCA
beamformer in synthesisable Verilog, integrate it with the single-channel pipeline and test the complete multichannel
system on the board against the Python reference. Work on the board directly (simulation only to debug a specific
problem), with separate agents for the existing FPGA design and board, the PCA fixed-point design, the RTL and the UART
data flow. Explain why we verify on the physical board and not only in Vivado simulation.
\end{prompt}

\begin{prompt}{A24}{4 Nov 2026}{\Anshu}
Before declaring the FPGA implementation successful, verify that the complete system really extracts the fetal ECG
correctly; correct extraction matters more than bit-exact agreement with Python. Use parallel agents to audit the PCA
algorithm, the fetal heart-rate detection and the Python--FPGA comparison, and validate the corrected system on real
four-channel recordings. Report F1, sensitivity and PPV on the benchmark, and fix the maternal false-alarm metric: the
report says 0\,\% but the CSV has nonzero values, so check whether maternal beats are detected as fetal beats.
\end{prompt}

\begin{prompt}{A25}{4--6 Nov 2026}{\Vivek}
Explain the nonzero LSB differences between Python and FPGA on five recordings (filter initial state or alignment), then
evaluate the design on the board: resource use, timing and throughput, latency and UART performance, with optimisation
where needed. Re-run the six Phase~4 cases on the board over COM5 with hardware mode mandatory (abort rather than fall
back to software), save the FPGA output for every sample separately from the Python reference, and report the maximum
and mean error in LSB, the exact-match percentage and the sample and packet counts. Then provide the exact script used
(with its SHA-256), show where the FPGA output array comes from, give the raw UART captures, explain what the timing
figures measure, and explain why an earlier log showed large errors while this run shows none.
\end{prompt}

\begin{prompt}{A26}{4 Nov 2026}{\Astitva}
Carry out a complete technical audit of Phases 1--5 without changing anything: list every file needed to verify the
reports, algorithms, RTL, data set results, UART communication and board tests independently (sources, testbenches,
constraints, Vivado projects, bitstreams, CSVs, timing and utilisation reports, board logs), separating essential from
optional evidence and flagging claims that lack evidence. Then collect these files unmodified into a ZIP with a manifest
of paths, sizes and SHA-256 hashes, listing any file that is missing.
\end{prompt}

\begin{prompt}{A27}{6 Nov 2026}{\Naman}
Expand the data set validation without changing the frozen algorithm, RTL, filter coefficients or bitstream. Audit the
FECGSYNDB coverage: find exactly which subject, SNR, layout (l1--l5), condition and repetition combinations have been
evaluated, separating floating-point, bit-accurate Python, RTL simulation and board tests, and produce a coverage matrix.
Then validate the frozen algorithm on the NIFEA clinical database (inspection, signal preparation, run, FPGA check) and
state the limits of that evaluation.
\end{prompt}

\begin{prompt}{A28}{6 Nov 2026}{\Astitva}
The Spartan-7 board is connected. Describe the current synthesis and timing-closure bottlenecks, then profile the RTL
datapath, find the critical-path latency and precision bottlenecks, and optimise the hardware to maximise causal fetal
QRS detection accuracy on the board.
\end{prompt}

% ---------------------------------------------------------------------------------------------
\section{Part B: Claude Code (Anthropic), 7 to 14 November 2026}
Model: Claude Opus~5.5.

\subsection{Causal algorithm on real recordings}

\begin{prompt}{B1}{7 Nov 2026}{\Anshu}
Audit the repository architecture, documentation and the Graphify knowledge graph, give a structured status of the
milestones, baseline metrics and bottlenecks across software and hardware, and confirm that the Spartan-7 board is
detected on the serial port. Commit the baseline, then audit and refine the extraction and detection algorithms to
maximise fetal QRS accuracy across all FECGSYNDB recordings and on the FPGA. Enforce strict physical causality
throughout: no non-causal filtering, no future-sample look-ahead and no ground-truth leakage. The PhysioNet/CinC
Challenge 2013 data is now in the data directory: extend the validation to these real recordings and redesign the
pipeline with causal filtering and online spatial beamforming.
\end{prompt}

\begin{prompt}{B2}{8 Nov 2026}{\Anshu}
Pause the CinC evaluation and report F1, sensitivity, PPV and FHR MAE on the other two benchmark data sets and on
NIFEADB. The real recordings score much lower than the synthetic ones: analyse the physiological and noise causes,
propose remedies, and implement them to improve robustness on clinical recordings.
\end{prompt}

\begin{prompt}{B3}{9 Nov 2026}{\Anshu}
Summarise the status of all algorithm variants. Evaluate the 1D convolutional U-Net on these data sets as a baseline
against our classical methods. Outline a strategy to improve the classical, hardware-oriented algorithms under real
clinical noise and changing maternal morphology, then implement it, preferring FPGA-friendly techniques over neural
networks: verify convergence and numerical stability in Python first, and move to RTL once the F1 targets are met.
\end{prompt}

\subsection{FPGA-centred design: RTL and PC split}

\begin{prompt}{B4}{9--10 Nov 2026}{\Astitva}
The Spartan-7 must be the real-time processing engine: the PC only streams the digitised samples, while channel
filtering, maternal suppression, QRS detection and FHR run on the chip; redesign the pipeline for this. Confirm whether
the bit-accurate Python models and the Verilog RTL are both implemented and verified, compare the floating-point
prototype with the fixed-point RTL (precision, latency, resources), and explain why we validate in software before RTL.
Evaluate a PC--FPGA split that leaves the host only the minimum (for example the initial calibration or covariance
decomposition); then commit and implement the RTL + PC split, run the multi-record validation, and compare the result
with all earlier baselines (F1, sensitivity, PPV, MAE, resources).
\end{prompt}

\begin{prompt}{B5}{13--14 Nov 2026}{\Astitva}
Summarise the milestones with the key metrics (accuracy, latency, resources). Explain why block RAM use on the
Spartan-7 is high and propose memory reductions, and give the exact hardware throughput and delay for a 1-second
(250-sample) input. Then implement the planned block RAM reductions.
\end{prompt}

\begin{prompt}{B6}{13--14 Nov 2026}{\Naman}
Write a remediation checklist before changing the architecture, and build a real-time web demonstration: upload
multi-lead recordings (PhysioNet WFDB \texttt{.dat}+\texttt{.hea} or 4-channel CSV), require the physical board on COM5,
stream at the real-time rate of 250\,Hz, show the live fetal heart rate, the UART connection and the processing stages,
and overlay the reference beats with accuracy figures when labels exist. Specify where the outputs are stored and explain
the steps the page performs when a recording is loaded.
\end{prompt}

\begin{prompt}{B7}{14 Nov 2026}{\Anshu}
Accuracy keeps improving over longer runs: analyse whether the adaptive thresholds and covariance estimates need a
settling time, and what a longer calibration window or more adaptation would do to the convergence and the steady-state
accuracy.
\end{prompt}

\subsection{Submission preparation, v6.1 and v7}

\begin{prompt}{B8}{14 Nov 2026}{\Naman}
Prepare the repository for submission: commit, restructure the folders and remove temporary build files, caches and
orphaned logs while keeping all source code, testbenches and validated reports (remove the duplicate deep-learning
environment \path{code/venv} and the superseded MicroBlaze v5 Vivado build tree but keep its bitstreams; keep the full
git history). Push, merge the release branch into main with the repository anonymous and the AI acknowledgements kept for
the report, and delete the merged feature branch. Propose a fair four-way division of the work over the whole project,
from preliminary exploration to hardware verification, so that each member's contribution is clear and defensible.
\end{prompt}

\begin{prompt}{B9}{14 Nov 2026}{\Naman}
Set up \path{report/final_report} from the IEEE/CVPR template of the earlier reports with all results, comparison tables,
observations and the contribution matrix, following the instructor's rules (pasted: one zip \texttt{fpr-<team \#>.zip}
with code, data and PDF; all new text in blue; all AI prompts in a separate file; AI acknowledgement mandatory) and with
an honest AI disclosure. Show which testbenches to simulate and which signals and time ranges to capture for Vivado
waveform figures. Suggest further analyses or verification sections, carry them out (except item~9), and rewrite
\path{README.md} with the architecture, installation and the steps to reproduce the software and board runs. Keep the AI
disclosure in the report short and the full prompt file in the archive.
\end{prompt}

\begin{prompt}{B10}{14 Nov 2026}{\Anshu}
Explain the main theoretical and physiological limits on accuracy in a purely causal single-pass design (moving dipole
axes, maternal QRS overlap, muscle noise). Implement the three proposed improvements on the causal baseline (v6.1), and
design a variant v7 with a bounded look-ahead window, future samples within a clinically acceptable buffer, to measure
the gain from non-causal spatial filtering and post-processing.
\end{prompt}

\begin{prompt}{B11}{14 Nov 2026}{\Astitva\ \Vivek}
Translate v6.1 into synthesisable RTL. Check whether the v7 look-ahead can run on the FPGA itself by buffering short
blocks in block RAM, so that the board outputs the results without continuous PC streaming, and extend the testbenches to
benchmark both versions. Find out why the current test run takes so long. Do not finish the report until v7 is verified:
stream the recordings in sub-blocks (a block latency of 15--20\,s is clinically acceptable), implement v7, synthesise it
and verify it on the board.
\end{prompt}

\begin{prompt}{B12}{14 Nov 2026}{\Naman}
Update the web demonstration to work with the v7 pipeline and show its intermediate results, give the commands to stop
old server instances, and suggest which recordings to load for the demo. Then identify the data set with the lowest
scores and explain why, and apply the final revisions to the report: remove obsolete sections, document v7 and its
results, and summarise the main additions.
\end{prompt}

\end{document}
